\SourceLeaf{409}

\[
\begin{aligned}
\operatorname{tang}\tfrac12 U
 &=\dfrac{\begin{gathered}y^{\frac18}\sin\dfrac12u-y^{\frac98}\sin\dfrac32u\\[3pt]+y^{\frac{25}8}\sin\dfrac52u-y^{\frac{49}8}\sin\dfrac72u+\mathord{.}\,\mathord{.}\end{gathered}}{\begin{gathered}y^{\frac18}\cos\dfrac12u+y^{\frac98}\cos\dfrac32u\\[3pt]+y^{\frac{25}8}\cos\dfrac52u+y^{\frac{49}8}\cos\dfrac72u+\mathord{.}\,\mathord{.}\end{gathered}}\\[14pt]
\operatorname{tang} U'
 &=\dfrac{\begin{gathered}y^{\frac14}\sin u-y^{\frac94}\sin3u\\[3pt]+y^{\frac{25}4}\sin5u-y^{\frac{49}4}\sin7u+\mathord{.}\,\mathord{.}\end{gathered}}{\begin{gathered}y^{\frac14}\cos u+y^{\frac94}\cos3u\\[3pt]+y^{\frac{25}4}\cos5u+y^{\frac{49}4}\cos7u+\mathord{.}\,\mathord{.}\end{gathered}}\\[14pt]
\operatorname{tang}2U''
 &=\dfrac{\begin{gathered}y^{\frac12}\sin2u-y^{\frac92}\sin6u\\[3pt]+y^{\frac{25}2}\sin10u-y^{\frac{49}2}\sin14u+\mathord{.}\,\mathord{.}\end{gathered}}{\begin{gathered}y^{\frac12}\cos2u+y^{\frac92}\cos6u\\[3pt]+y^{\frac{25}2}\cos10u+y^{\frac{49}2}\cos14u+\mathord{.}\,\mathord{.}\end{gathered}}
\end{aligned}
\]
[wird. Für den Satz, dass die Reihen $P$, $Q$ den Functionalgleichungen genügen,
welche den besprochenen Algorithmus bestimmen, und welche oben in Art.~22
zusammengestellt sind, hat \textsc{Gauss} ausser dem Beweise, der sich auf die
Verwandlung jener Reihen in unendliche Producte stützt und der in der unten
folgenden Abhandlung `hundert Theoreme über die neue Transscendente' enthalten
ist, wahrscheinlich auch noch einen andern Beweis geführt, wie er sich leicht
aus den oben in Art.~16 gemachten Andeutungen ergibt.]

\AGMArticle{26}
[Bezeichnen wir, abweichend von der in den vorhergehenden Artikeln befolgten
Weise, die ersten Derivirten der Functionen $P(y,\allowbreak \eta)$, $Q(y,\allowbreak \eta)$, $R(y,\allowbreak \eta)$,
$S(y,\allowbreak \eta)$ nach der Grösse $\log\eta$ als unabhängige Veränderliche mit
$P'$, $Q'$, $R'$, $S'$ und die zweiten Derivirten nach derselben Grösse mit
$P''$, $Q''$, $R''$, $S''$, ferner die ersten Derivirten der Functionen
$py$, $qy$, $ry$ nach $\log y$ als unabhängige Veränderliche mit $p'$, $q'$, $r'$,
so folgt aus den für $a$, $b$, $c$, $\alpha$, $\beta$, $\gamma$, $\delta$ gefundenen
Differentialgleichungen:
\[
\begin{aligned}\dfrac{qr'-rq'}{qrp^4}&=\dfrac{pr'-rp'}{rpq^4}
 \\[7pt]&=\dfrac{qp'-pq'}{pqr^4}\\[7pt]&=\tfrac14\end{aligned}
\]
\[
\begin{aligned}\dfrac{PS'-SP'}{ppQR}&=\dfrac{QS'-SQ'}{qqPR}
 \\[7pt]&=\dfrac{RS'-SR'}{rrPQ}\\[7pt]&=\dfrac{QR'-RQ'}{ppPS}
 \\[7pt]&=\dfrac{PR'-RP'}{qqQS}\\[7pt]&=\dfrac{QP'-PQ'}{rrRS}\\[7pt]&=1\end{aligned}
\]
\[
\begin{aligned}-4\dfrac{Q'}Q&=(ry)^2\mathbin{.}\dfrac{S(yy,\eta\eta)}{Q(yy,\eta\eta)}\\[6pt]&\quad{}+
(ryy)^2\mathbin{.}\dfrac{S(y^4,\eta^4)}{Q(y^4,\eta^4)}\\[6pt]&\quad{}+
(ry^4)^2\mathbin{.}\dfrac{S(y^8,\eta^8)}{Q(y^8,\eta^8)}+\mathord{.}\,\mathord{.}\end{aligned}
\]
\[
\begin{aligned}
\dfrac{P'P'-PP''}{PP}+\tfrac14pprr\dfrac{RR}{PP}
 &=\dfrac{Q'Q'-QQ''}{QQ}-\tfrac14pprr\dfrac{SS}{QQ}
\\[8pt]&=\dfrac{R'R'-RR''}{RR}+\tfrac14pprr\dfrac{PP}{RR}
 \\[8pt]&=\dfrac{S'S'-SS''}{SS}-\tfrac14pprr\dfrac{QQ}{SS}
\\[8pt]&=-\dfrac{q'}q\\[8pt]&=\tfrac18\{(ry)^4+2(ryy)^4+4(ry^4)^4+8(ry^8)^4+\mathord{.}\,\mathord{.}\}
\end{aligned}
\]

\SourceLeaf{410}

\[
\dfrac{\mathrm d\mathrm dP}{(\mathrm d\log\eta)^2}
 =\dfrac{\mathrm dP}{\mathrm d\log y},
\]
welch letzterer Gleichung jede der Functionen $P$, $Q$, $R$, $S$ genügt.
Die vorhergehende mehrfache Gleichung zwischen den ersten und zweiten
Derivirten der Functionen nach der Grösse $\log\eta$ findet sich, soweit sie
sich auf $P$ und $S$ bezieht, in \textsc{Gauss} handschriftlichem Nachlasse an
einer von den übrigen Untersuchungen dieser Functionen getrennten Stelle.
Es sind dort $P$, $S$, $p$, $r$ in einer für diese specielle Entwickelung etwas
bequemeren Form als die hier benutzte durch ihre Reihen definirt, und es
heisst dann,] \textit{so wird}
\[
\begin{aligned}
P'P'-PP''&=-\dfrac{\mathrm dp(yy)}{\mathrm d\log y}\mathbin{.}P(yy,\eta\eta)\\[5pt]&\quad{}-
\dfrac{\mathrm dr(yy)}{\mathrm d\log y}\mathbin{.}R(yy,\eta\eta)\\[9pt]
S'S'-SS''&=+\dfrac{\mathrm dr(yy)}{\mathrm d\log y}\mathbin{.}P(yy,\eta\eta)\\[5pt]&\quad{}-
\dfrac{\mathrm dp(yy)}{\mathrm d\log y}\mathbin{.}R(yy,\eta\eta)
\end{aligned}
\]
\[
\begin{aligned}
PP&=+p(yy)\mathbin{.}P(yy,\eta\eta)+r(yy)\mathbin{.}R(yy,\eta\eta)\\
SS&=-r(yy)\mathbin{.}P(yy,\eta\eta)+p(yy)\mathbin{.}R(yy,\eta\eta)
\end{aligned}
\]
\textit{hier ist noch zu bemerken (wovon jedoch der Beweis tiefer liegt)}
\[
\begin{aligned}&p(yy)\mathbin{.}\dfrac{\mathrm dr(yy)}{\mathrm d\log y}
-r(yy)\mathbin{.}\dfrac{\mathrm dp(yy)}{\mathrm d\log y}
\\[6pt]&=\tfrac12p(yy)\mathbin{.}r(yy)\mathbin{.}\{p(yy)^4-r(yy)^4\}\end{aligned}
\]
\textit{also}
\[
\begin{aligned}&\dfrac{P'P'-PP''}{PP}-\dfrac{S'S'-SS''}{SS}
\\[6pt]&=-\tfrac12\left(\dfrac{PP}{SS}+\dfrac{SS}{PP}\right)
\mathbin{.}p(yy)\mathbin{.}r(yy)\mathbin{.}\{p(yy)^2-r(yy)^2\}\end{aligned}
\]
\textit{noch findet man}
\[
\begin{aligned}
PS'-SP'={}&\tfrac12\left\{\begin{gathered}y^{\frac18}+3y^{\frac98}\\[3pt]-5y^{\frac{25}8}
 -7y^{\frac{49}8}+9y^{\frac{81}8}+\mathord{.}-\mathord{.}-\mathord{.}\end{gathered}\right\}\times\\[10pt]
&\times\left\{\begin{gathered}y^{\frac18}(\eta^{\frac12}+\eta^{-\frac12})
\\[3pt]-y^{\frac98}(\eta^{\frac32}+\eta^{-\frac32})
\\[3pt]-y^{\frac{25}8}(\eta^{\frac52}+\eta^{-\frac52})
 +\mathord{.}+\mathord{.}-\mathord{.}-\end{gathered}\right\}
\end{aligned}
\]
\textit{Das Quadrat des zweiten Factors im andern Theile der vorstehenden Gleichung wird}
\[
=r\mathbin{.}Q(y,\eta\eta)+q\mathbin{.}R(y,\eta\eta)
\]
\textit{Der erste Factor wird}
\[
=\{p(yy)^2+r(yy)^2\}{}^4\!\surd \tfrac12
 \{p(yy)^2-r(yy)^2\}p(yy)ryy)\ ?
\]
\textit{Zusammen wird, reductis reducendis}


\SourceLeaf{411}

\[
\begin{aligned}
(PS'-SP')^2={}&\tfrac14\{p(yy)^2+r(yy)^2\}\\[6pt]&\quad\{p(yy)P(yy,\eta\eta)-r(yy)R(yy,\eta\eta)\}\times\\
&\times\{r(yy)P(yy,\eta\eta)+p(yy)R(yy,\eta\eta)\}\\[10pt]
={}&\tfrac14\{[p(yy)^2-r(yy)^2]PP+2p(yy)r(yy)SS\}\times\\
&\times\{2p(yy)r(yy)PP-[p(yy)^2-r(yy)^2]SS\}
\end{aligned}
\]
\Needspace{16\baselineskip}
\textit{Setzt man also}
\[
\eta=e^{i\varphi}
\]
\[
\dfrac{-iS}{P}\surd {\dfrac{p(yy)^2-r(yy)^2}{2p(yy)r(yy)}}=\sin\theta
\]
\textit{so wird}
\[
\mathrm d\varphi=
\dfrac{2\,\mathrm d\theta}
 {\surd {\left(\begin{gathered}{}[p(yy)^2-r(yy)^2]\cos\theta^2\\[3pt]{}+
[p(yy)^2+r(yy)^2]\sin\theta^2\end{gathered}\right)}}
\]
\[
\cos\theta=\dfrac RP\surd {\dfrac{p(yy)^2+r(yy)^2}{2p(yy)r(yy)}}
\]
\begin{center}\rule{0.23\textwidth}{0.3pt}\end{center}
